\documentclass[12pt]{article}

\usepackage{amsmath, amssymb, amsthm}
\usepackage{enumitem}
\usepackage{geometry}
\usepackage{fancyhdr}
\usepackage{hyperref}

% --- Page setup ---
\geometry{margin=1in}
\pagestyle{fancy}
\fancyhf{}
\rhead{MATH 421 -- 003, Fall 2025}
\lhead{Homework 5}
\cfoot{\thepage}
\renewcommand{\headrulewidth}{0.4pt}

% --- Title info ---
\title{Homework 5}
\author{Alejandro De La Torre \\ MATH 421 -- The Theory of Single Variable Calculus \\ Section 003 \\ Instructor: Grace Work}
\date{October 2025}

\begin{document}
\maketitle

\section*{Collaboration Statement}
Work compared with Griffin Hailman, with whom brief advice and discussion were exchanged on some problems.

% ---------------- Problem 1 ----------------
\section*{Problem 1 (Ross Chapter 2, 7.4)}
Give examples of

\begin{enumerate}[label=\textbf{(\alph*)}]
  \item A sequence $(x_n)$ of irrational numbers having a limit $\lim x_n$ that is a rational number.
  \item A sequence $(r_n)$ of rational numbers having a limit $\lim r_n$ that is an irrational number.
\end{enumerate}

\paragraph{Discussion}
\textit{Foreword}: I realize now that I did not read and all I had to do is provide examples that did not need proofs. Still, despite 2 hours of sunk cost, I am going to keep this (probably incomplete) proof here because I feel mildly proud of it and hope that, in the case where my subsequent work lacks, I hope this justifies some additional points (pretty please). \smallskip

    This problem asks us to construct two specific kinds of sequences using what we know about rational and irrational numbers.
    
    A \emph{sequence} is simply a function whose domain is the natural numbers (or a tail of them), so each integer $n$ corresponds to a real value $s_n$.  We say that a sequence \emph{converges} to a real number $s$ if for every $\varepsilon>0$ there exists some $N$ such that $|s_n-s|<\varepsilon$ whenever $n>N$.
    
    Recall that a number is \emph{rational} if it can be expressed as $\tfrac{m}{n}$ for integers $m,n$ with $n\neq0$; otherwise it is \emph{irrational}.  The set of rational numbers is closed under addition, subtraction, and multiplication.  From this it follows that:
    the product of a nonzero rational number and an irrational number is irrational, and the sum of a rational and an irrational number is irrational.
    
    For \textit{part (a)} we want every term of the sequence to be irrational, but the sequence should approach a rational number as $n$ grows.  One natural idea is to start with a rational target, say $1$, and then \emph{perturb} it by a small irrational piece that tends to $0$.  This ensures each term stays irrational, while the irrational part vanishes in the limit.
    
    For \textit{part (b)} we reverse the situation: every term should be rational, but the sequence should approach an irrational limit.  Here we can use the \emph{density of the rationals in the reals}, which guarantees that rationals can approximate any real number---even an irrational one---as closely as we like.


\paragraph{Proof for (a)}
Let
\[
x_n = 1 + \frac{\sqrt{2}}{n}, \qquad n\in\mathbb{N}.
\]
Since $\sqrt{2}$ is irrational, $\tfrac{\sqrt{2}}{n}$ must also be irrational: if it were rational, then multiplying by the rational number $n$ would make $\sqrt{2}$ rational, which is impossible.  Furthermore, adding a rational number (here $1$) to an irrational number produces an irrational number.  Hence every $x_n$ is irrational.

Because $\lim\tfrac{1}{n}=0$, we have
\[
\lim_{n\to\infty}x_n
  = \lim_{n\to\infty}\!\Bigl(1+\frac{\sqrt{2}}{n}\Bigr)
  = 1 + \sqrt{2}\!\cdot\!\lim_{n\to\infty}\!\frac{1}{n}
  = 1+0=1.
\]
Therefore the sequence $(x_n)$ consists of irrational numbers and converges to the rational number $1$. In other words, $(x_n)$ is irrational for all $n$, and $\displaystyle\lim x_n=1$.

\paragraph{Proof for (b)}
We want to construct a sequence of rational numbers that converges to an irrational number.  
Since $\sqrt{2}$ is irrational, it is a natural choice for our limit.  
The idea is to take rational numbers that approximate $\sqrt{2}$ more and more closely.

\smallskip
Let
\[
r_n = 1.4,\, 1.41,\, 1.414,\, 1.4142,\, \dots
\]
denote the sequence obtained by truncating the decimal expansion of $\sqrt{2}$ after $n$ digits.  
Formally, each $r_n$ is a rational number because it has only finitely many decimal places.  
As we add more digits of $\sqrt{2}$, these rational numbers get arbitrarily close to $\sqrt{2}$.

\smallskip
To express this idea symbolically, we can write
\[
|r_n - \sqrt{2}| < \frac{1}{10^n}.
\]
Since $\frac{1}{10^n} \to 0$ as $n \to \infty$, it follows that
\[
\lim_{n \to \infty} r_n = \sqrt{2}.
\]

\smallskip
Therefore, $(r_n)$ is a sequence of rational numbers whose limit is the irrational number $\sqrt{2}$.  
In summary:
\[
r_n \in \mathbb{Q} \text{ for all } n, 
\qquad \text{and} \qquad
\lim_{n \to \infty} r_n = \sqrt{2}.
\]

% ---------------- Problem 2 ----------------
\section*{Problem 2 (Ross Chapter 2, 8.1c)}
Prove the following:
\[
\lim_{n \to \infty} \frac{2n - 1}{3n + 2} = \frac{2}{3}.
\]

\paragraph{Discussion}
Intuitively, the highest powers of \(n\) dominate both the numerator and the denominator:
\((2n-1)\) behaves like \(2n\) and \((3n+2)\) behaves like \(3n\), so the ratio should
approach \(2n/3n = 2/3\). A tidy way to formalize this with limit laws is to divide
numerator and denominator by \(n\), which turns the “lower order” terms into \(1/n\)-type
pieces that go to \(0\). If we want an explicit convergence argument, we can
also compute the difference from \(2/3\) exactly and bound it by \(\varepsilon\).

\paragraph{Proof (via limit laws)}
Write
\[
\frac{2n-1}{3n+2}
= \frac{\,\displaystyle 2-\frac1n\,}{\,\displaystyle 3+\frac2n\,}.
\]
Since \(\lim_{n\to\infty}\tfrac1n=\lim_{n\to\infty}\tfrac2n=0\), we have
\[
\lim_{n\to\infty}\Bigl(2-\tfrac1n\Bigr)=2
\quad\text{and}\quad
\lim_{n\to\infty}\Bigl(3+\tfrac2n\Bigr)=3\neq 0.
\]
By the quotient law for limits,
\[
\lim_{n\to\infty}\frac{2n-1}{3n+2}
=\frac{\;\lim(2-\tfrac1n)\;}{\;\lim(3+\tfrac2n)\;}
=\frac{2}{3}.
\]

\paragraph{Proof (via Convergence)}
For every \(n\),
\[
\left|\frac{2n-1}{3n+2}-\frac23\right|
= \left|\frac{(2n-1)(3)-(3n+2)(2)}{(3n+2)\,3}\right|
= \frac{7}{9n+6}.
\]
Given \(\varepsilon>0\), choose
\[
N>\frac{1}{9}\left(\frac{7}{\varepsilon}-6\right).
\]
Then for all \(n\ge N\),
\[
\left|\frac{2n-1}{3n+2}-\frac23\right|=\frac{7}{9n+6}
\le \frac{7}{9N+6}<\varepsilon.
\]
Hence $\lim_{n\to\infty}\frac{2n-1}{3n+2}=\frac{2}{3}$.

% ---------------- Problem 3 ----------------
\section*{Problem 3 (Ross Chapter 2, 8.4)}
Let $(t_n)$ be a bounded sequence, i.e., there exists $M$ such that $|t_n| \le M$ for all $n$, and let $(s_n)$ be a sequence such that $\lim s_n = 0$. Prove $\lim (s_n t_n) = 0$.

\paragraph{Discussion}
So we are given that $|t_n| \leq M \space \forall \space n \in \mathbb{N}$ such that $t_n$ is never going to be greater than M. 
We are asked to prove that if one sequence $(s_n)$ tends to zero and another sequence $(t_n)$
remains bounded, then their product $(s_n t_n)$ also tends to zero. 

The intuition is that $(t_n)$ cannot grow without bound—there is a fixed number $M$ such that 
$|t_n| \le M$ for all $n$. As $(s_n)$ becomes smaller and smaller (approaching $0$), multiplying 
it by any bounded number cannot “blow it up.” In fact, the inequality 
\[
|s_n t_n| = |s_n||t_n| \le M|s_n|
\]
captures this relationship precisely. Because $|s_n|$ can be made as small as we wish, $M|s_n|$
can also be made arbitrarily small. Therefore, $(s_n t_n)$ must converge to $0$.

\paragraph{Proof}
Let $\varepsilon > 0$ be given. Since $\lim s_n = 0$, there exists $N$ such that
\[
|s_n| < \frac{\varepsilon}{M} \quad \text{for all } n > N.
\]
Because $(t_n)$ is bounded, we have $|t_n| \le M$ for all $n$. Then, for $n > N$,
\[
|s_n t_n| = |s_n||t_n| \le M|s_n| < M \cdot \frac{\varepsilon}{M} = \varepsilon.
\]
Thus, by the definition of convergence, $\lim_{n \to \infty} (s_n t_n) = 0$.\qed


% ---------------- Problem 4 ----------------
\section*{Problem 4 (Ross Chapter 2, 8.5)}
\begin{enumerate}[label=\textbf{(\alph*)}]
  \item Consider three sequences $(a_n)$, $(b_n)$, and $(s_n)$ such that $a_n \le s_n \le b_n$ for all $n$ and $\lim a_n = \lim b_n = s$. Prove $\lim s_n = s$. \\
  \textit{This is called the “squeeze lemma.”}
  \item Suppose $(s_n)$ and $(t_n)$ are sequences such that $|s_n| \le t_n$ for all $n$ and $\lim t_n = 0$. Prove $\lim s_n = 0$.
\end{enumerate}

\paragraph{Solution to (a)}
\subparagraph{Discussion}
The idea is: for the same index $n$, each $s_n$ sits between $a_n$ and $b_n$.  
As $a_n$ and $b_n$ both get arbitrarily close to $s$, the “gap’’ \,$[a_n,b_n]$ shrinks to the single
point $s$. There is no room left for $s_n$ to go anywhere else, so $s_n$ must also approach $s$.

\subparagraph{Proof}
Let $\varepsilon>0$. Since $\lim a_n=s$, there is $N_1$ such that $|a_n-s|<\varepsilon$ for all $n>N_1$.
Similarly, since $\lim b_n=s$, there is $N_2$ such that $|b_n-s|<\varepsilon$ for all $n>N_2$.
With $N=\max\{N_1,N_2\}$, for all $n>N$ we have
\[
s-\varepsilon<a_n\le s_n\le b_n<s+\varepsilon,
\]
so $|s_n-s|<\varepsilon$. Hence, by the definition of convergence, $\lim s_n=s$.
\qed

\paragraph{Solution to (b)}
\subparagraph{Discussion}
The hypothesis $|s_n|\le t_n$ forces $t_n\ge0$. If $t_n\to0$, then eventually $t_n$ is smaller than any
prescribed $\varepsilon>0$. The squeeze $-t_n\le s_n\le t_n$ then pins $s_n$ within $\varepsilon$ of $0$,
so $s_n\to0$. (Formally, this is just the squeeze lemma with $a_n=-t_n$, $b_n=t_n$, $s=0$.)

\subparagraph{Proof}
Let $\varepsilon>0$. Because $\lim t_n=0$, there exists $N$ such that $n>N$ implies $t_n<\varepsilon$.
For such $n$ we have
\[
|s_n-0|=|s_n|\le t_n<\varepsilon.
\]
Therefore $s_n\to0$ by the definition of convergence.  \qed


% ---------------- Problem 5 ----------------
\section*{Problem 5 (Ross Chapter 2, 8.8b)}
Prove the following:
\[
\lim_{n \to \infty} [\sqrt{n^2 + n} - n] = \frac{1}{2}.
\]
\paragraph{Discussion}
We want to prove $\displaystyle\lim_{n\to\infty}\bigl(\sqrt{n^2+n}-n\bigr)=\tfrac12$ by the
$\varepsilon$–$N$ definition. The obstacle is the square root.
A standard trick in these situations is to \emph{rationalize} the
expression we must control; i.e., multiply and divide by a suitable
conjugate so that the square root disappears from the numerator.
Here the target limit is $\tfrac12$, so we study
\[
\Bigl|\sqrt{n^2+n}-n-\frac12\Bigr|.
\]
If we set $a=\sqrt{n^2+n}$ and $b=n+\tfrac12$, then
$a-b=\dfrac{a^2-b^2}{a+b}$, which removes the square root from the
numerator. This produces an expression whose size is $\sim\! \dfrac{1}{n}$
because the denominator grows linearly with $n$. Once we bound the
difference by a constant multiple of $\tfrac{1}{n}$, the
$\varepsilon$–$N$ argument is routine.

\paragraph{Proof}
Let $\varepsilon>0$ be given. Write
\[
\sqrt{n^2+n}-n-\frac12
=\frac{(\sqrt{n^2+n}-n-\tfrac12)\bigl(\sqrt{n^2+n}+n+\tfrac12\bigr)}
      {\sqrt{n^2+n}+n+\tfrac12}
=\frac{(n^2+n)-(n+\tfrac12)^2}{\sqrt{n^2+n}+n+\tfrac12}.
\]
Since $(n+\tfrac12)^2=n^2+n+\tfrac14$, the numerator simplifies to
$-\tfrac14$. Hence
\[
\Bigl|\sqrt{n^2+n}-n-\frac12\Bigr|
=\frac{1}{4\bigl(\sqrt{n^2+n}+\,n+\tfrac12\bigr)}.
\]
Because $\sqrt{n^2+n}\ge n$ for all $n\ge 1$, we have
\[
4\bigl(\sqrt{n^2+n}+\,n+\tfrac12\bigr)
\;\ge\;4\bigl(n+n+\tfrac12\bigr)
\;=\;8n+2,
\]
and therefore
\[
\Bigl|\sqrt{n^2+n}-n-\frac12\Bigr|
\;\le\;\frac{1}{8n+2}
\;<\;\frac{1}{8n}\qquad(n\ge1).
\]
Choose $N\in\mathbb{N}$ with $N>\dfrac{1}{8\varepsilon}$.
Then for all $n\ge N$,
\[
\Bigl|\sqrt{n^2+n}-n-\frac12\Bigr|
\;<\;\frac{1}{8n}
\;\le\;\frac{1}{8N}
\;<\;\varepsilon.
\]
By the definition of convergence, $\displaystyle\lim_{n\to\infty}
\bigl(\sqrt{n^2+n}-n\bigr)=\tfrac12$.
\qed
% ---------------- Problem 6 ----------------
\section*{Problem 6 (Ross Chapter 2, 8.9)}
Let $(s_n)$ be a sequence that converges.

\begin{enumerate}[label=\textbf{(\alph*)}]
  \item Show that if $s_n \ge a$ for all but finitely many $n$, then $\lim s_n \ge a$.
  \item Show that if $s_n \le b$ for all but finitely many $n$, then $\lim s_n \le b$.
  \item Conclude that if all but finitely many $s_n$ belong to $[a,b]$, then $\lim s_n$ belongs to $[a,b]$.
\end{enumerate}

\paragraph{Solution to (a)}
\subparagraph{Discussion}
If a convergent sequence $(s_n)$ is eventually bounded \emph{below} by $a$, then its tail never dips under $a$. 
Intuitively, once $s_n$ is within any $\varepsilon>0$ of its limit $s$, all far–out terms lie in $(s-\varepsilon,s+\varepsilon)$.
If $s<a$, we can pick $\varepsilon$ so small that $(s-\varepsilon,s+\varepsilon)$ lies strictly below $a$, which contradicts
the fact that all sufficiently large $s_n$ are $\ge a$. Hence $s\ge a$.

\subparagraph{Proof}
Let $(s_n)$ converge to $s$ and assume $s_n\ge a$ for all but finitely many $n$.
Suppose, for a contradiction, that $s<a$. Set $\varepsilon:=\dfrac{a-s}{2}>0$.
Because $s_n\to s$, there exists $N_1$ such that $|s_n-s|<\varepsilon$ for all $n>N_1$, hence
\[
s_n<s+\varepsilon \;=\; \frac{s+a}{2} \;<\; a \qquad(n>N_1).
\]
Also, since only finitely many $n$ violate $s_n\ge a$, there exists $N_2$ with $s_n\ge a$ for all $n>N_2$.
Taking $N=\max\{N_1,N_2\}$ yields a contradiction for $n>N$. Therefore $s\ge a$. \qed

\bigskip

\paragraph{Solution to (b)}
\subparagraph{Discussion}
This is the upper-bound analogue of part (a). If $s_n\le b$ eventually but the limit were $s>b$,
then for a small $\varepsilon$ the tail $(s-\varepsilon,s+\varepsilon)$ would lie strictly above $b$, 
contradicting $s_n\le b$ for all large $n$. Hence $s\le b$.

\subparagraph{Proof}
Let $(s_n)\to s$ and assume $s_n\le b$ for all but finitely many $n$.
Suppose, for a contradiction, that $s>b$. Set $\varepsilon:=\dfrac{s-b}{2}>0$.
Then there exists $N_1$ with $|s_n-s|<\varepsilon$ for all $n>N_1$, so
\[
s_n> s-\varepsilon \;=\; \frac{s+b}{2} \;>\; b \qquad(n>N_1).
\]
Since only finitely many $n$ violate $s_n\le b$, there exists $N_2$ such that $s_n\le b$ for all $n>N_2$.
Taking $N=\max\{N_1,N_2\}$ gives a contradiction for $n>N$. Hence $s\le b$. \qed

\bigskip

\paragraph{Solution to (c)}
\subparagraph{Discussion}
“All but finitely many $s_n$ lie in $[a,b]$” means that beyond some index, $s_n\ge a$ and $s_n\le b$ simultaneously.
Parts (a) and (b) then force $a\le \lim s_n \le b$, i.e., the limit belongs to $[a,b]$.

\subparagraph{Proof}
If all but finitely many $s_n$ lie in $[a,b]$, then for some $N$ we have $a\le s_n\le b$ for all $n>N$.
By (a) we get $\lim s_n\ge a$, and by (b) we get $\lim s_n\le b$. Therefore $\lim s_n\in[a,b]$. \qed

\end{document}